% ============ 目录（手排，不用 \tableofcontents，因为节头是自定义宏） ============
\subhdr{本册目录}

\begin{minipage}{\linewidth}
\small
\begin{tabular}{@{}>{\raggedright\arraybackslash}p{0.08\linewidth}>{\raggedright\arraybackslash}p{0.30\linewidth}>{\raggedright\arraybackslash}p{0.56\linewidth}@{}}
\S0 & 讲义档案 & 12 份课件的坐标、三种标记、读法、页数交代 \\
\S1 & 语音与字母 & 33 个字母分三类；元音五对；重音与 о 弱化；硬软辅音；清浊与两条清化；三个调型；五个基本问答 \\
\S2 & 名词 & кто／что；三个性与 -ь 的 38 词；复数三表合一；24 组特殊复数；只有单／复数 \\
\S3 & 人称与物主代词 & 八个主语；ты／Вы；四列式物主代词；先判性再选形；Меня́ зову́т 与问候 \\
\S4 & 形容词 & 位置与 како́й 四形式；五族词尾总表；拼写规则七字母；13 个颜色词；三步与四个坑 \\
\S5 & 指示代词与 нра́виться & э́тот 四形式；са́мый；нра́виться ＋ 第三格；что／потому́ что；ИК-5；课文 ПАРК СИНХА́Й \\
\S6 & 动词 & 两个变位组；-ся 动词；-л 型过去时；бу́ду ＋ 不定式；三时态对照；по-ру́сски 与副词 \\
\S7 & 数词 & 1--1000；读长数字；价格（юа́нь 三形）；时间（час 三形）；序数词；星期与月份 \\
\S8 & 复习卡 & 四张速查表：名词／词尾／动词／数词 \\
\S9 & 易错点 & 十二条，每条都能指回某个词的某一步 \\
\S10 & 缺口清单 & 12 条：课件跳了哪一步，本讲义补在哪 \\
\S11 & 三句话总结 & 判性数、一条规律多个位置、能回指页码 \\
\S12 & 页码对照表 & 12 份课件 → 本讲义哪一节 \\
\end{tabular}
\end{minipage}

\vspace{2pt}

{\small\color{inkgray}读法建议：第一次顺着 \S1 到 \S7 读；考前只读 \S8 与 \S9；做练习卡住时先查 \S8 的表，再回对应节的「关键一步」。}
